\chapter{Determinants and Trace}
\label{ch:i08-determinants-and-trace}

The determinant compresses invertibility, oriented volume scaling and alternating multilinearity into one scalar. Trace measures the sum of diagonal action and is the simplest similarity invariant. Together they are the first examples of basis-independent quantities extracted from a matrix representation.

\section*{Learning goals}
After this chapter, the reader should be able to:
\begin{itemize}
\item compute determinants efficiently using row operations, triangular form, or cofactor expansion;
\item track the effect of row swaps, row scalings, and row additions on the determinant;
\item use the determinant to decide invertibility and linear independence for square matrices;
\item prove and apply multiplicativity \(\det(AB)=\det(A)\det(B)\);
\item compute traces and use cyclicity \(\operatorname{tr}(AB)=\operatorname{tr}(BA)\) in concrete identities;
\item interpret determinant and trace as basis-independent invariants of a linear operator;
\end{itemize}
\section*{Conceptual roadmap}
The chapter is organized around a repeated pattern:
\[
\boxed{\text{definition}\;\longrightarrow\;\text{structural theorem}\;\longrightarrow\;\text{coordinate calculation}\;\longrightarrow\;\text{application}.}
\]
Whenever possible, we separate the invariant mathematical object from a chosen coordinate representation.

\section{Alternating multilinear characterization}
The determinant is the unique alternating multilinear function of the columns with $\det I=1$.

\section{Permutation formula}
For $A=(a_{ij})$, $\det A=\sum_{\sigma\in S_n}\operatorname{sgn}(\sigma)\prod_i a_{i,\sigma(i)}$.

\section{Row operations}
Swapping rows changes sign, scaling a row scales the determinant, and adding a multiple of one row to another leaves it unchanged.

\section{Invertibility}
A square matrix is invertible exactly when its determinant is nonzero.


% BEGIN VOL01-EXPANSION I08-example-02
\begin{example}[A determinant hypothesis that matters]\label{ex:i08-expand-02}
If \(A\) and \(B\) are square matrices, then
\[
\det(AB)=\det A\,\det B.
\]
Consequently \(AB\) is invertible exactly when both factors are invertible:
a zero determinant in either factor forces \(\det(AB)=0\), while two nonzero
determinants give a nonzero product.

This gives a scalar certificate for a structural fact about composition of
linear isomorphisms.
\end{example}
% END VOL01-EXPANSION I08-example-02
\section{Geometric meaning}
Over $\mathbb R$, $|\det A|$ is the factor by which the associated linear map scales $n$-dimensional volume; the sign records orientation.

\section{Trace}
$\operatorname{tr}(A)=\sum_i a_{ii}$ is linear and satisfies $\operatorname{tr}(AB)=\operatorname{tr}(BA)$, hence is invariant under similarity.


% BEGIN VOL01-EXPANSION I08-example-01
\begin{example}[Determinant as an invertibility check]\label{ex:i08-expand-01}
For
\[
A=\begin{pmatrix}2&1\\3&2\end{pmatrix},
\qquad
\det A=4-3=1.
\]
Thus \(A\) is invertible. Solving
\(A(x,y)^T=(5,8)^T\) gives \((x,y)=(2,1)\), and substitution checks the
solution.

The determinant does not produce the solution by itself; it tells us that a
unique solution must exist for every right-hand side.
\end{example}
% END VOL01-EXPANSION I08-example-01

% BEGIN VOL01-EXPANSION I08-example-03
\begin{example}[Trace survives change of basis]\label{ex:i08-expand-03}
Let \(B=P^{-1}AP\). Using cyclicity of trace,
\[
\operatorname{tr}(B)
=\operatorname{tr}(P^{-1}AP)
=\operatorname{tr}(APP^{-1})
=\operatorname{tr}(A).
\]
Hence trace is attached to the operator, not to one particular matrix
representation.

A practical check follows: if two matrices are claimed to be similar but have
different traces, the claim is impossible.
\end{example}
% END VOL01-EXPANSION I08-example-03
\section{Core structural results}

\begin{theorem}[Multiplicativity of determinant]\label{thm:i08-01}
For square matrices $A,B$, $\det(AB)=\det(A)\det(B)$.
\begin{proof}
For fixed $B$, the map sending the columns of $A$ to $\det(AB)$ is alternating multilinear and has normalization $\det(B)$ at $A=I$; uniqueness gives the formula.
\end{proof}
\end{theorem}

\begin{theorem}[Cyclic trace identity]\label{thm:i08-02}
For compatible square matrices, $\operatorname{tr}(AB)=\operatorname{tr}(BA)$.
\begin{proof}
Expand the diagonal entries: $\operatorname{tr}(AB)=\sum_{i,j}a_{ij}b_{ji}$ and interchange the dummy indices.
\end{proof}
\end{theorem}

\section{Worked examples}

\begin{example}[Triangular matrices]\label{ex:i08-worked-01}
The determinant of a triangular matrix is the product of its diagonal entries, and its trace is their sum.
\end{example}

\begin{example}[Area scaling]\label{ex:i08-worked-02}
The matrix $\begin{pmatrix}2&1\\0&3\end{pmatrix}$ scales signed area by $6$.
\end{example}

\begin{example}[Shear]\label{ex:i08-worked-03}
A shear $\begin{pmatrix}1&t\\0&1\end{pmatrix}$ has determinant $1$ although it changes angles.
\end{example}

\section{Diagnostic checklist}
Before moving on, it is useful to distinguish four levels of understanding:
\begin{enumerate}
\item recognize the definitions in unfamiliar examples;
\item prove the core structural statements without coordinates when possible;
\item translate the same statement into a matrix or coordinate computation;
\item know which hypotheses are essential and produce a counterexample when one is removed.
\end{enumerate}

% BEGIN VOLUME01 SOLVED DOSSIERS
\section{Solved dossiers}
These dossiers are longer worked problems. They complement the shorter exercise--hint--solution triads by combining several ideas in one argument.

\begin{problem}[Determinant by elimination]\label{prob:i08-dossier-01}
Compute $\det\begin{pmatrix}1&2&3\\0&1&4\\0&0&5\end{pmatrix}$ and explain why no expansion is needed.
\end{problem}
\begin{solution}
The matrix is upper triangular, so the determinant is the product of diagonal entries: $1\cdot1\cdot5=5$.
\end{solution}

\begin{problem}[Effect of a row replacement]\label{prob:i08-dossier-02}
What happens to the determinant when row $2$ is replaced by row $2+7$ row $1$?
\end{problem}
\begin{solution}
The determinant is unchanged. Multilinearity expands the new determinant into the old one plus $7$ times a determinant with two equal rows, and the latter is zero by alternation.
\end{solution}

\begin{problem}[Area scaling]\label{prob:i08-dossier-03}
What signed area factor is produced by $A=\begin{pmatrix}2&1\\1&3\end{pmatrix}$?
\end{problem}
\begin{solution}
$\det A=6-1=5$. Thus oriented area is multiplied by $5$; ordinary area is multiplied by $|\det A|=5$.
\end{solution}

\begin{problem}[Invertibility test]\label{prob:i08-dossier-04}
Why does $\det A=0$ imply the columns are dependent?
\end{problem}
\begin{solution}
The determinant is alternating multilinear and vanishes on dependent columns. Conversely, if the columns were independent they would form a basis and the associated linear map would be invertible, forcing nonzero determinant.
\end{solution}

\begin{problem}[Trace cyclicity]\label{prob:i08-dossier-05}
Verify directly that $\operatorname{tr}(AB)=\operatorname{tr}(BA)$.
\end{problem}
\begin{solution}
$\operatorname{tr}(AB)=\sum_i\sum_j a_{ij}b_{ji}$. Interchanging the dummy indices gives $\sum_j\sum_i b_{ji}a_{ij}=\operatorname{tr}(BA)$.
\end{solution}

\begin{problem}[Similarity invariants]\label{prob:i08-dossier-06}
Show determinant and trace are invariant under $A\mapsto PAP^{-1}$.
\end{problem}
\begin{solution}
Multiplicativity gives $\det(PAP^{-1})=\det P\,\det A\,\det P^{-1}=\det A$. Cyclicity gives $\operatorname{tr}(PAP^{-1})=\operatorname{tr}(AP^{-1}P)=\operatorname{tr}A$.
\end{solution}

\begin{problem}[Determinant of a reflection]\label{prob:i08-dossier-07}
Find the determinant of reflection across a hyperplane in $\mathbb R^n$.
\end{problem}
\begin{solution}
The reflection fixes the $(n-1)$-dimensional hyperplane and negates one normal direction. In an adapted basis the matrix is $\operatorname{diag}(1,\dots,1,-1)$, so the determinant is $-1$.
\end{solution}

\begin{problem}[Determinant of a shear]\label{prob:i08-dossier-08}
Why does $\begin{pmatrix}1&t\\0&1\end{pmatrix}$ preserve area for every $t$?
\end{problem}
\begin{solution}
Its determinant is $1$. Geometrically a shear slides parallel slices without changing base or height, consistent with area preservation.
\end{solution}

\begin{problem}[Trace of rank one]\label{prob:i08-dossier-09}
For column vectors $u,v$, compute $\operatorname{tr}(uv^T)$.
\end{problem}
\begin{solution}
The diagonal entries are $u_iv_i$, so the trace is $\sum_i u_iv_i=v^Tu$.
\end{solution}

\begin{problem}[Block triangular determinant]\label{prob:i08-dossier-10}
If $M=\begin{pmatrix}A&B\\0&C\end{pmatrix}$ with square diagonal blocks, what is $\det M$?
\end{problem}
\begin{solution}
The determinant is $\det A\,\det C$. This follows by expansion or by reducing to triangular block operations; the off-diagonal block $B$ does not affect the product of diagonal block determinants.
\end{solution}

\begin{problem}[Orientation test]\label{prob:i08-dossier-11}
How can determinant distinguish the two connected components of $GL_n(\mathbb R)$?
\end{problem}
\begin{solution}
An invertible real matrix has determinant either positive or negative. Continuous paths in $GL_n$ cannot cross determinant $0$, so the sign cannot change along a path. Positive and negative determinants correspond to the two orientation classes.
\end{solution}

\begin{problem}[Determinant synthesis]\label{prob:i08-dossier-12}
Why is the determinant better viewed as a property of an endomorphism than of one matrix?
\end{problem}
\begin{solution}
Changing basis replaces the matrix by a similar one, and determinant is similarity invariant. Thus $\det T$ is well-defined independently of coordinates; a matrix merely computes it.
\end{solution}

% END VOLUME01 SOLVED DOSSIERS

\section{Exercises with complete solutions}

\begin{exercise}\label{exr:i08-01}
Compute $\det\begin{pmatrix}a&b\\c&d\end{pmatrix}$.
\end{exercise}
\begin{hint}
Use row operations to reach triangular form, but record which operations change the determinant and how.
\end{hint}
\begin{solution}
The permutation formula gives $ad-bc$.
\end{solution}

\begin{exercise}\label{exr:i08-02}
Show that a matrix with two equal columns has determinant zero.
\end{exercise}
\begin{hint}
Reduce to triangular form while recording how each row operation changes the determinant. Expand along a row or column with the most zeros rather than using a full permutation formula.
\end{hint}
\begin{solution}
Alternation changes the determinant sign when the equal columns are swapped, but the matrix is unchanged, so the determinant equals its negative.
\end{solution}

\begin{exercise}\label{exr:i08-03}
Why is $\det A\ne0$ equivalent to invertibility?
\end{exercise}
\begin{hint}
Place the vectors as columns and ask when the determinant vanishes; that is exactly the dependence condition in the square case. A square matrix is singular exactly when its determinant vanishes; connect that to a nontrivial kernel.
\end{hint}
\begin{solution}
Nonzero determinant implies full column independence; equivalently the associated endomorphism has zero kernel and hence is invertible.
\end{solution}

\begin{exercise}\label{exr:i08-04}
Compute the determinant of a diagonal matrix.
\end{exercise}
\begin{hint}
Reduce to triangular form while recording how each row operation changes the determinant. For \(\det(AB)\), first handle the invertible case or use the effect of elementary matrices.
\end{hint}
\begin{solution}
Only the identity permutation contributes, so the determinant is the product of diagonal entries.
\end{solution}

\begin{exercise}\label{exr:i08-05}
Prove $\operatorname{tr}(PAP^{-1})=\operatorname{tr}(A)$.
\end{exercise}
\begin{hint}
Write the diagonal entries of \(AB\) as a double sum and interchange the summation indices. The determinant of a triangular matrix is the product of its diagonal entries.
\end{hint}
\begin{solution}
Use cyclicity: $\operatorname{tr}(PAP^{-1})=\operatorname{tr}(AP^{-1}P)=\operatorname{tr}(A)$.
\end{solution}

\begin{exercise}\label{exr:i08-06}
What determinant does a reflection of $\mathbb R^n$ across a hyperplane have?
\end{exercise}
\begin{hint}
First determine \(\lambda\) from \(\det(A-\lambda I)=0\), then solve \((A-\lambda I)v=0\) for the eigenspace. For trace, write the diagonal entries of \(AB\) and interchange the summation indices to compare with \(BA\).
\end{hint}
\begin{solution}
It has eigenvalues $1$ on the hyperplane and $-1$ on a normal direction, so determinant $-1$.
\end{solution}

\begin{exercise}\label{exr:i08-07}
Compute the trace of a rank-one matrix $uv^T$.
\end{exercise}
\begin{hint}
Similarity preserves determinant and trace because the factors \(P\) and \(P^{-1}\) cancel inside these invariants.
\end{hint}
\begin{solution}
The trace is $v^Tu=\sum_i u_iv_i$.
\end{solution}

\begin{exercise}\label{exr:i08-08}
Show that determinant and trace are basis-independent for an endomorphism.
\end{exercise}
\begin{hint}
Place the vectors as columns and ask when the determinant vanishes; that is exactly the dependence condition in the square case. Use determinant as a quick consistency check for an alleged inverse: \(\det(A^{-1})=1/\det(A)\).
\end{hint}
\begin{solution}
Two matrix representations are similar; determinant and trace are invariant under similarity.
\end{solution}


% BEGIN VOL01-EXPANSION I08-exercises-01
\section{Graded supplementary exercises}
The following exercises extend the original exercise set without replacing it. They are grouped by mathematical role and increase in synthesis.

\subsection*{Standard computations}

\begin{exercise}[Two-by-two determinant]\label{exr:i08-expand-01}
Compute \(\det\begin{pmatrix}3&1\\2&4\end{pmatrix}\).
\end{exercise}
\begin{hint}
Use \(ad-bc\).
\end{hint}
\begin{solution}
\(12-2=10\).
\end{solution}

\begin{exercise}[Triangular determinant]\label{exr:i08-expand-02}
Compute the determinant of an upper-triangular matrix with diagonal \(2,-1,5\).
\end{exercise}
\begin{hint}
Multiply diagonal entries.
\end{hint}
\begin{solution}
\(-10\).
\end{solution}

\begin{exercise}[Trace]\label{exr:i08-expand-03}
Compute the trace of \(\begin{pmatrix}1&2\\3&-4\end{pmatrix}\).
\end{exercise}
\begin{hint}
Add diagonal entries.
\end{hint}
\begin{solution}
\(-3\).
\end{solution}

\begin{exercise}[Scaling determinant]\label{exr:i08-expand-04}
For \(A\in M_3\), what is \(\det(2A)\)?
\end{exercise}
\begin{hint}
Scaling all three rows by \(2\) multiplies determinant by \(2^3\).
\end{hint}
\begin{solution}
\(\det(2A)=8\det A\).
\end{solution}

\begin{exercise}[Invertibility]\label{exr:i08-expand-05}
Can a matrix with determinant \(0\) be invertible?
\end{exercise}
\begin{hint}
Use \(\det(AA^{-1})=1\).
\end{hint}
\begin{solution}
No; invertibility would force \(\det A\det A^{-1}=1\).
\end{solution}

\subsection*{Proofs}

\begin{exercise}[Alternating property]\label{exr:i08-expand-06}
Prove a matrix with two equal rows has determinant zero.
\end{exercise}
\begin{hint}
Swap the equal rows.
\end{hint}
\begin{solution}
Alternation changes the sign but leaves the matrix unchanged, so \(\det A=-\det A\); over characteristic not \(2\), this gives zero, and the determinant axioms give zero generally.
\end{solution}

\begin{exercise}[Multiplicativity consequence]\label{exr:i08-expand-07}
Prove if \(AB\) is invertible then both square factors are invertible.
\end{exercise}
\begin{hint}
Use determinants.
\end{hint}
\begin{solution}
\(\det(AB)=\det A\det B\neq0\), so both determinants are nonzero.
\end{solution}

\begin{exercise}[Trace cyclicity]\label{exr:i08-expand-08}
Prove \(\operatorname{tr}(AB)=\operatorname{tr}(BA)\).
\end{exercise}
\begin{hint}
Expand diagonal entries as double sums.
\end{hint}
\begin{solution}
Both equal \(\sum_{i,j}a_{ij}b_{ji}\).
\end{solution}

\begin{exercise}[Similarity invariance]\label{exr:i08-expand-09}
Prove determinant is invariant under similarity.
\end{exercise}
\begin{hint}
Compute \(\det(P^{-1}AP)\).
\end{hint}
\begin{solution}
\(\det(P^{-1})\det(A)\det(P)=\det A\).
\end{solution}

\subsection*{Counterexamples}

\begin{exercise}[Counterexample: trace product]\label{exr:i08-expand-10}
Is \(\operatorname{tr}(AB)=\operatorname{tr}A\,\operatorname{tr}B\) always true?
\end{exercise}
\begin{hint}
Test identity matrices.
\end{hint}
\begin{solution}
No. For \(A=B=I_2\), the left side is \(2\), the right side \(4\).
\end{solution}

\begin{exercise}[Counterexample: determinant sum]\label{exr:i08-expand-11}
Is \(\det(A+B)=\det A+\det B\) generally true?
\end{exercise}
\begin{hint}
Test \(A=B=I_2\).
\end{hint}
\begin{solution}
No: \(\det(2I_2)=4\) but \(\det I_2+\det I_2=2\).
\end{solution}

\begin{exercise}[Counterexample: same trace]\label{exr:i08-expand-12}
Do equal traces imply similarity?
\end{exercise}
\begin{hint}
Compare \(I_2\) and a nontrivial Jordan block with eigenvalue \(1\).
\end{hint}
\begin{solution}
No; both have trace \(2\), but one is diagonalizable and the other is not.
\end{solution}

\subsection*{Applications}

\begin{exercise}[Application: area scaling]\label{exr:i08-expand-13}
What oriented area factor does \(\begin{pmatrix}2&0\\0&3\end{pmatrix}\) produce?
\end{exercise}
\begin{hint}
Use the determinant.
\end{hint}
\begin{solution}
The factor is \(6\).
\end{solution}

\begin{exercise}[Application: orientation]\label{exr:i08-expand-14}
What does a negative determinant say about an invertible real planar map?
\end{exercise}
\begin{hint}
Interpret the sign of determinant.
\end{hint}
\begin{solution}
It reverses orientation.
\end{solution}

\subsection*{Challenges}

\begin{exercise}[Challenge: rank-one update]\label{exr:i08-expand-15}
Let \(A=\begin{pmatrix}1&t\\0&1\end{pmatrix}\). Compute determinant and trace for all \(t\).
\end{exercise}
\begin{hint}
Use triangular formulas.
\end{hint}
\begin{solution}
\(\det A=1\), \(\operatorname{tr}A=2\), independent of \(t\); these invariants do not detect the shear parameter.
\end{solution}

\begin{exercise}[Challenge: commutator trace]\label{exr:i08-expand-16}
Show \(\operatorname{tr}(AB-BA)=0\).
\end{exercise}
\begin{hint}
Use cyclicity of trace.
\end{hint}
\begin{solution}
\(\operatorname{tr}(AB)-\operatorname{tr}(BA)=0\).
\end{solution}
% END VOL01-EXPANSION I08-exercises-01
\section*{Chapter summary}
The determinant compresses invertibility, oriented volume scaling and alternating multilinearity into one scalar. Trace measures the sum of diagonal action and is the simplest similarity invariant. Together they are the first examples of basis-independent quantities extracted from a matrix representation.
The important habit is to move fluently between invariant statements and concrete coordinates while remembering which conclusions depend on finite dimensionality, the scalar field, or the inner product.
